\section{Worked computational workflow and result interpretation}
\label{sec:worked-workflow}

This chapter connects the preceding equations to the order of operations in
a restricted closed-shell Gamma calculation.  It is not a molecule-specific
input tutorial; it follows which object is built at each step and which
numerical representation owns it.

\subsection{Step 0: define the numerical problem}

Consider an RSDF source backend, the \texttt{multigrid\_xc} target integrator,
and r2SCAN.  Before the calculation, fix
\begin{itemize}
  \item coordinates, lattice, units, charge, and spin;
  \item primary basis, GTH pseudopotential, and auxiliary basis;
  \item Gamma or an explicit k-point list and weights;
  \item source cutoff/mesh and RSDF source/metric $\omega$ and meshes;
  \item target multigrid mesh; and
  \item HF, gradient, and Z-vector tolerances.
\end{itemize}
Recording both source and target meshes defines a reproducible problem more
completely than saying only that the cutoff was 500 Ha.

\subsection{Step 1: source HF}

RSDF builds $J[D]$ and $K[D]$ in the iteration
\begin{equation}
  F_{\mathrm{HF}}=h+J[D]-\tfrac12K[D].
\end{equation}
After convergence, retain $D^{\mathrm{HF}}$, MO coefficients, occupations,
orbital energies, and the source DF provider.  There is no r2SCAN target
optimization in this step.  If switching the target integrator from uniform
to multigrid changes this source state, the ownership implementation is
incorrect.

\subsection{Step 2: fixed-density target energy}

Evaluate the target functional at the same $D^{\mathrm{HF}}$.  In M1,
\begin{equation}
  E_{\mathrm{target}}
  =\Tr[D^{\mathrm{HF}}h]+E_J^{\mathrm{RSDF}}[D^{\mathrm{HF}}]
   +E_{\mathrm{xc}}^{\mathrm{MG}}[D^{\mathrm{HF}}]+E_{NN},
\end{equation}
and
\begin{equation}
  F_{\mathrm{target}}
  =h+J_{\mathrm{RSDF}}[D^{\mathrm{HF}}]
    +V_{\mathrm{xc}}^{\mathrm{MG}}[D^{\mathrm{HF}}].
\end{equation}
No SCF is run with this Fock matrix.  The energy workspace retains the target
Fock and multigrid tasks for reuse by the gradient.

\subsection{Step 3: nonstationary right-hand side}

Project the target Fock into the source HF VO block,
\begin{equation}
  g_{ai}=2(C_v^\dagger F_{\mathrm{target}}C_o)_{ai}.
\end{equation}
If $g=0$, the target energy happens to be stationary in the HF orbitals and
there is no orbital-response contribution.  In a general DC-DFT calculation,
$g\ne0$ and becomes the Z-vector right-hand side.

\subsection{Step 4: source HF Z-vector}

For small real dimension $m$, a dense $A_R$ provides a reference solution.
Large calculations use matrix-free AX with RSDF transition-density J/K and
MINRES.  The report should contain
\begin{itemize}
  \item dense or matrix-free solver selection;
  \item real response dimension and iteration count;
  \item minimum orbital gap and preconditioner;
  \item final scaled residual $\eta_z$; and
  \item total and mean-iteration solver wall times.
\end{itemize}
A small iteration count alone is not an accuracy certificate.  The explicit
residual and small-system dense parity are the completion evidence.

\subsection{Step 5: explicit target gradient}

At fixed target density, evaluate
\begin{itemize}
  \item selected one-electron $h^{Ax}$ and GTH PP derivatives;
  \item the RSDF target $dJ/dR$;
  \item the multigrid target $dE_{\mathrm{xc}}/dR$;
  \item the target energy-weighted-density contraction with $dS/dR$; and
  \item the nuclear--nuclear derivative.
\end{itemize}
Because this is M1, RSDF owns target $dJ/dR$ and multigrid contributes XC
only.  In M2, combined multigrid $dJ+dE_{\mathrm{xc}}$ is used and the RSDF
target-$dJ$ route is disabled.

\subsection{Step 6: source response gradient}

For each $Ax$, build the one-electron, overlap, and source-RSDF two-electron
nuclear source $B^{Ax}$ and accumulate $z^TB^{Ax}$.  Neither full $B$ nor full
$C^{Ax}$ needs to remain resident.  Nuclear-component batches, source/metric
adjoint contractions, and reciprocal streaming leave only an $(N_A,3)$
scalar array.

\subsection{Step 7: final assembly}

Equation~\eqref{eq:final-gradient-assembly} gives
\begin{equation}
  \mathbf g
  =\mathbf g_{\mathrm{explicit,target}}
   +\mathbf g_{\mathrm{Pulay}}
   +\mathbf g_{\mathrm{source,response}}
   +\mathbf g_{NN}.
\end{equation}
The public output is a gradient, not a force, and is reported in Ha/Bohr.

\subsection{Shape and ownership trace}

\begin{center}
\begin{longtable}{>{\raggedright\arraybackslash}p{0.22\textwidth}
                  >{\raggedright\arraybackslash}p{0.25\textwidth}
                  >{\raggedright\arraybackslash}p{0.20\textwidth}
                  >{\raggedright\arraybackslash}p{0.22\textwidth}}
\toprule
Object & Representative shape & Owner & Next consumer \\
\midrule
$D^{\mathrm{HF}}$ & $(N_k,N_{AO},N_{AO})$ & source HF & target energy, AX \\
$F_{\mathrm{target}}$ & same & target integrator + DF & VO right-hand side \\
$g_{vo}$ & per-k $(N_v,N_o)$ & common orbital layer & Z solver \\
$z_{vo}$ & per-k $(N_v,N_o)$ & source HF Hessian & nuclear response \\
Target XC gradient & $(N_A,3)$ & multigrid & final assembly \\
RSDF response & $(N_A,3)$ & RSDF adjoint & final assembly \\
Final gradient & $(N_A,3)$ & common assembly & report/FD \\
\bottomrule
\end{longtable}
\end{center}

\subsection{Validation order}

A completed large calculation is not sufficient validation.  A useful order
is
\begin{enumerate}
  \item primitive matrix/response parity on a small distorted geometry;
  \item Gamma versus explicit one-k-point shape parity;
  \item sampled-k conjugation and $1/N_k$ checks;
  \item selected Cartesian full-energy finite differences;
  \item a complete $3N_A$ small-system FD gate;
  \item separate convergence of source mesh, target mesh, auxiliary basis,
        and k-mesh; and
  \item larger-system timing and memory measurements.
\end{enumerate}

Analytic--FD agreement establishes that a chosen finite representation was
differentiated correctly.  Mesh convergence separately establishes that the
finite representation is close enough to the desired physical limit.

\subsection{Reading timing and memory reports}

Decompose end-to-end time into source HF, target energy, target explicit
gradient, Z-vector, one-electron/PP, DF response, and final reporting.  High
CPU use in a GPU calculation may reflect host integral generation, task
construction, or packing and does not by itself imply a final CPU fallback.
Use ownership metadata together with stage timing.

Peak VRAM must also be assigned to a stage.  OOM is an allocation failure,
whereas \texttt{CUDA\_ERROR\_ILLEGAL\_ADDRESS} is an invalid access by a
launched kernel.  CUDA is asynchronous, so the next CuPy operation in a
traceback may not be the first failing operation.  A large-grid error should
be localized with blocking diagnostics before it is hidden by lowering a
cutoff; the physical representation and a kernel limitation are separate
issues.

\subsection{Further reading}

The periodic Gaussian software framework is described in
\href{https://doi.org/10.1002/wcms.1340}{the PySCF framework paper}.  The
range-separated three-center construction used by RSDF is described by Ye
and Berkelbach,
\href{https://doi.org/10.1063/5.0046617}{J. Chem. Phys. 154, 131104 (2021)}.
The SCAN and r2SCAN definitions are given in
\href{https://doi.org/10.1103/PhysRevLett.115.036402}{Phys. Rev. Lett. 115,
036402 (2015)} and
\href{https://doi.org/10.1021/acs.jpclett.0c02405}{J. Phys. Chem. Lett. 11,
8208--8215 (2020)}, respectively.  Implementation-specific ownership and
validation mappings are maintained in
\texttt{docs/implementation/Equation\_Code\_Map.md}.
